\clearpage\thispagestyle{empty}
\pdfbookmark[0]{Quellen zum Pallas-Bericht vom Juli 1803}{pallas180307-sources}
\begin{center}\textbf{Quellen zum Pallas-Bericht vom Juli 1803}\end{center}
\noindent Diese KI-Übertragung in heutiges Deutsch folgt dem Auszug aus dem Pallas-Bericht von Franz Xaver von Zach in \href{https://gdz.sub.uni-goettingen.de/id/PPN236007467}{Gauß’ \textit{Werken}, Band~VI} (Göttingen, 1874), S.~245. Der Auszug entspricht dem letzten Absatz und der letzten Tabelle des Beitrags VII im Juliheft der \textit{Monatlichen Correspondenz}, Band~VIII (1803), \href{https://api.digitale-sammlungen.de/iiif/image/v2/bsb10538600_00096/full/1450,/0/default.jpg}{S.~90--93}, auf \href{https://api.digitale-sammlungen.de/iiif/image/v2/bsb10538600_00099/full/1450,/0/default.jpg}{S.~93}. Die Tabelle der \textit{Werke} gibt die mittleren Bremer Zeiten an und bezeichnet die Unterschiede mit R—B. Die Beobachtungen stammen von Heinrich Wilhelm Olbers, der Vergleich mit den Elementen VI von Carl Friedrich Gauß.
\par\smallskip\noindent Übertragung und Satz wurden mit KI erstellt. Die Digitalisate der Zeitschrift stammen von der Bayerischen Staatsbibliothek.
\clearpage
